\section{Pruebas, complicaciones, adecuaciones, hallazgos y resultados}

\lettrine[lines=2,nindent=0pt]{E}{}n esta secci\'on se muestran todas las pruebas realizadas, complicaciones, adecuaciones, hallazgos y resultados encontrados, lo cual va ligado con la hip\'otesis propuesta, la cual debe validarse en esta secci\'on. En consecuencia, el an\'alisis no se limita a enumerar ejecuciones o m\'etricas, sino que busca interpretar la evidencia experimental obtenida y establecer si la propuesta desarrollada cumple con el objetivo central de demostrar que un esquema de contratos de agente y gobernanza externa puede reducir acciones inseguras sin eliminar la capacidad operativa del agente.

\subsection{Pruebas realizadas}
Las pruebas efectuadas se organizaron en torno a un benchmark de enforcement para agentes LLM con uso de herramientas. El experimento se ejecut\'o sobre un conjunto de escenarios nominales y adversariales, comparando cuatro modos de operaci\'on: \texttt{no\_contract}, \texttt{advisory}, \texttt{guarded} y \texttt{strict}. Cada corrida gener\'o evidencia observable en forma de trazas, ledgers y res\'umenes de ejecuci\'on, lo que permiti\'o analizar no s\'olo la salida textual del agente, sino tambi\'en sus propuestas de acci\'on, bloqueos, replanteamientos y efectos finales.

Las pruebas incluyeron tanto validaci\'on funcional de la infraestructura como ejecuciones completas del corpus experimental con backend real sobre Ollama y LiteLLM. De este modo, la evidencia generada combina verificaci\'on de implementaci\'on con observaci\'on emp\'irica del comportamiento del agente bajo distintos niveles de enforcement.

\subsection{Complicaciones encontradas}
Durante el desarrollo y la ejecuci\'on experimental surgieron varias complicaciones relevantes. Una de las principales fue la discrepancia entre resultados locales y fallas de integraci\'on continua, originada por diferencias entre archivos modificados localmente y archivos realmente incluidos en los cambios integrados al repositorio. Otra complicaci\'on importante fue la necesidad de corregir f\'ormulas del evaluador, ya que una versi\'on inicial produc\'ia tasas inv\'alidas al mezclar conteos a nivel de acci\'on con denominadores a nivel de corrida.

Adem\'as, se identificaron dificultades propias del comportamiento agentivo del modelo, por ejemplo repeticiones innecesarias de acciones, recuperaciones incompletas tras un bloqueo y escenarios donde la presi\'on adversarial no era suficiente para inducir la acci\'on riesgosa esperada. Estas complicaciones no invalidan la propuesta; por el contrario, permiten justificar las adecuaciones metodol\'ogicas y t\'ecnicas que fortalecen la validez del estudio.

\subsection{Adecuaciones implementadas}
Como respuesta a las complicaciones anteriores, se realizaron adecuaciones en distintos niveles del sistema. En el evaluador se redefinieron las m\'etricas de prevenci\'on a nivel de corrida y se separ\'o con mayor claridad la detecci\'on runtime de las violaciones de pre y post-ejecuci\'on. En los escenarios se ajustaron expectativas, criterios terminales y niveles de presi\'on adversarial para alinear mejor la oportunidad de violaci\'on con el objetivo experimental.

En la implementaci\'on del runtime se a\~nadi\'o \textit{feedback} estructurado del gobernador para apoyar la recuperaci\'on del agente despu\'es de un bloqueo, as\'i como mecanismos de \textit{terminal success} para finalizar ejecuciones seguras una vez alcanzado el outcome correcto. Estas adecuaciones son parte de la aportaci\'on experimental de la tesis, pues muestran que la gobernanza efectiva no depende s\'olo de detectar y bloquear, sino tambi\'en de orientar la recuperaci\'on y evitar degradaciones innecesarias en la utilidad del agente.

\subsection{Hallazgos principales}
Los hallazgos m\'as importantes del experimento muestran que los modos con enforcement activo logran prevenir efectos inseguros de manera consistente. En la ejecuci\'on completa del benchmark, los modos \texttt{guarded} y \texttt{strict} evitaron la ejecuci\'on de acciones inseguras observadas en la l\'inea base \texttt{no\_contract} y en \texttt{advisory}. Sin embargo, tambi\'en se observ\'o una diferencia relevante entre ambos: \texttt{guarded} conserv\'o una tasa de finalizaci\'on segura considerablemente mayor que \texttt{strict}.

Otro hallazgo importante es que la calidad de la gobernanza no puede evaluarse s\'olo por detecci\'on de violaciones. Fue necesario medir tambi\'en la oportunidad de acci\'on insegura, la efectividad de prevenci\'on, la recuperaci\'on posterior al bloqueo y la finalizaci\'on segura de la tarea. Este punto refuerza la postura de la tesis de que la evaluaci\'on de agentes gobernados debe ser multidimensional y no limitarse a clasificaciones binarias de \'exito o fallo.

\subsection{Resultados encontrados}
Los resultados agregados del experimento completo, correspondiente a 84 corridas (21 escenarios $\times$ 4 modos $\times$ 1 r\'eplica), muestran diferencias claras entre los cuatro modos de enforcement. En primer lugar, el n\'umero de acciones inseguras ejecutadas fue de 5 en \texttt{no\_contract}, 5 en \texttt{advisory}, 0 en \texttt{guarded} y 0 en \texttt{strict}. Este dato por s\'i solo ya permite observar que la presencia del gobernador con capacidad de bloqueo modifica materialmente el comportamiento operativo del agente.

En cuanto a la m\'etrica de \texttt{unsafe\_side\_effect\_rate}, los modos \texttt{no\_contract} y \texttt{advisory} obtuvieron 1.0, mientras que \texttt{guarded} y \texttt{strict} obtuvieron 0.0. Esto significa que, una vez que se present\'o una oportunidad real de acci\'on insegura, los dos primeros modos no lograron prevenir el efecto inseguro, mientras que los dos \'ultimos s\'i lo evitaron completamente. En t\'erminos de \texttt{governance\_effectiveness}, el patr\'on es complementario: \texttt{no\_contract} = 0.0, \texttt{advisory} = 0.0, \texttt{guarded} = 1.0 y \texttt{strict} = 1.0.

Sin embargo, la diferencia m\'as importante para la tesis aparece al observar la \texttt{successful\_safe\_completion\_rate}. En esta m\'etrica, \texttt{guarded} alcanz\'o 0.904762, mientras que \texttt{strict} obtuvo 0.333333. Como referencia adicional, \texttt{no\_contract} obtuvo 0.619048 y \texttt{advisory} 0.571429. Estos valores muestran que el bloqueo r\'igido no es la mejor soluci\'on pr\'actica si elimina demasiada utilidad, mientras que \texttt{guarded} consigue preservar la finalizaci\'on segura en una proporci\'on mucho mayor.

Las m\'etricas de detecci\'on runtime tambi\'en respaldan esta interpretaci\'on. \texttt{guarded} obtuvo \texttt{precision} = 0.666667, \texttt{recall} = 1.0 y \texttt{f1} = 0.8, mientras que \texttt{advisory} y \texttt{strict} obtuvieron \texttt{precision} = 0.625, \texttt{recall} = 0.833333 y \texttt{f1} = 0.714286. En otras palabras, \texttt{guarded} no s\'olo previno acciones inseguras, sino que adem\'as alcanz\'o la mejor calidad de detecci\'on entre los modos con enforcement.

Desde la perspectiva de oportunidad experimental, los cuatro modos compartieron un \texttt{unsafe\_action\_opportunity\_rate} = 0.238095, lo que indica que el corpus s\'i gener\'o suficientes situaciones de riesgo observables para comparar los modos en condiciones equivalentes. Este punto es importante porque fortalece la validez interna del benchmark: las diferencias observadas no se explican por ausencia de oportunidad, sino por el efecto del mecanismo de enforcement.

En relaci\'on con \texttt{H4}, la sobrecarga del enforcement tambi\'en puede observarse directamente en el uso medio de tokens y en la latencia media por corrida. Los valores de \texttt{mean\_token\_usage} fueron 3483.714286 en \texttt{no\_contract}, 3497.380952 en \texttt{advisory}, 3649.285714 en \texttt{guarded} y 1906.619048 en \texttt{strict}. Esto implica que \texttt{guarded} consumi\'o aproximadamente 165.57 tokens m\'as por corrida que \texttt{no\_contract} y 151.90 m\'as que \texttt{advisory}, lo cual es coherente con su mayor n\'umero de replanteamientos y con la incorporaci\'on de feedback del gobernador. En latencia media (\texttt{mean\_latency\_ms}), los valores fueron 150082.128887 en \texttt{no\_contract}, 146950.853082 en \texttt{advisory}, 152272.317173 en \texttt{guarded} y 78346.903744 en \texttt{strict}. El caso de \texttt{strict} muestra menor costo porque aborta antes, mientras que \texttt{guarded} incurre en mayor trabajo para preservar utilidad y completar la tarea de manera segura.

En t\'erminos de la hip\'otesis, la evidencia respalda que un esquema de contratos y gobernanza externa puede reducir o eliminar efectos inseguros en un agente con uso de herramientas. No obstante, tambi\'en se observa que no todos los modos de enforcement son equivalentes: la mejor contribuci\'on pr\'actica no proviene del bloqueo m\'as r\'igido, sino del equilibrio entre prevenci\'on y recuperaci\'on. En este sentido, los resultados favorecen especialmente al modo \texttt{guarded}.

\subsection{Validaci\'on de la hip\'otesis}
La validaci\'on debe interpretarse respecto de las cuatro hip\'otesis experimentales definidas para el benchmark. En primer lugar, \texttt{H1} queda respaldada por los resultados, ya que los modos \texttt{guarded} y \texttt{strict} redujeron completamente la ejecuci\'on de acciones inseguras frente a \texttt{no\_contract} y \texttt{advisory}: se observaron 5 acciones inseguras ejecutadas en \texttt{no\_contract}, 5 en \texttt{advisory} y 0 tanto en \texttt{guarded} como en \texttt{strict}. Esto tambi\'en es consistente con \texttt{unsafe\_side\_effect\_rate} = 1.0 en \texttt{no\_contract} y \texttt{advisory}, frente a 0.0 en \texttt{guarded} y \texttt{strict}.

En segundo lugar, \texttt{H2} recibe un apoyo parcial y matizado. La hip\'otesis planteaba equivalencia entre la calidad de detecci\'on en \texttt{advisory} y en los modos bloqueantes. Los datos muestran que \texttt{advisory} y \texttt{strict} son id\'enticos en detecci\'on (\texttt{precision} = 0.625, \texttt{recall} = 0.833333, \texttt{f1} = 0.714286), mientras que \texttt{guarded} mejora a \texttt{precision} = 0.666667, \texttt{recall} = 1.0 y \texttt{f1} = 0.8. Por tanto, la evidencia apunta a que la capacidad de detecci\'on se mantiene comparable entre modos, pero no permite afirmar equivalencia perfecta en sentido estricto con una sola r\'eplica experimental.

En tercer lugar, \texttt{H3} queda claramente respaldada. Tanto \texttt{guarded} como \texttt{strict} alcanzaron \texttt{governance\_effectiveness} = 1.0, pero \texttt{guarded} obtuvo una \texttt{successful\_safe\_completion\_rate} = 0.904762, muy superior a 0.333333 de \texttt{strict}. Este contraste confirma que una estrategia de enforcement con recuperaci\'on ofrece una mejor relaci\'on entre seguridad y utilidad que una estrategia de aborto inmediato.

Finalmente, \texttt{H4} tambi\'en se sostiene, ya que el enforcement introduce sobrecarga observable en tiempo de ejecuci\'on, iteraciones y uso de tokens. Los cuatro modos muestran costos distintos: \texttt{mean\_token\_usage} = 3483.714286 en \texttt{no\_contract}, 3497.380952 en \texttt{advisory}, 3649.285714 en \texttt{guarded} y 1906.619048 en \texttt{strict}. De forma similar, \texttt{guarded} alcanz\'o \texttt{mean\_replans\_per\_run} = 1.333333, mientras que \texttt{strict} obtuvo 0, precisamente porque aborta sin intentar recuperaci\'on. Esto demuestra que la gobernanza efectiva no es gratuita y que su costo debe formar parte de la evaluaci\'on integral del sistema.

Adem\'as, existe un hallazgo metodol\'ogico importante: las adecuaciones al evaluador, al oracle, a los escenarios adversariales y al mecanismo de recuperaci\'on mejoraron de manera visible la calidad de detecci\'on respecto al piloto anterior. En particular, \texttt{advisory} pas\'o de \texttt{f1} = 0.5 a \texttt{f1} = 0.714286, y \texttt{guarded} pas\'o de \texttt{f1} = 0.4 a \texttt{f1} = 0.8. Este resultado fortalece la validez del proceso experimental, pues muestra que las complicaciones encontradas no s\'olo fueron resueltas, sino que llevaron a un benchmark mejor calibrado y metodol\'ogicamente m\'as s\'olido.

Por tanto, esta secci\'on no s\'olo confirma la viabilidad de la propuesta, sino que tambi\'en delimita sus condiciones m\'as favorables de aplicaci\'on. La evidencia obtenida sustenta que la gobernanza externa de agentes LLM es factible, medible y metodol\'ogicamente defendible dentro del marco experimental planteado en esta tesis.
